% ============================================================
%  Expanded Section: Fundamentals of Photodetection
%  Ge-on-Si Waveguide Photodiode Thesis
% ============================================================

\section[Fundamentals of Photodetection]{Fundamentals of Photodetection}
\label{sec:pd_fundamentals}

Photodetection begins with the transfer of optical energy to the semiconductor
lattice. Understanding that process at the physics level is a prerequisite for
rational device design, because every performance metric defined in
\cref{sec:pd_metrics} traces back to the rate at which photons are absorbed,
how the resulting carriers are transported, and what fraction is lost to
recombination before reaching the device terminals.

% ============================================================
\subsection[Optical Absorption in Semiconductors]{Optical Absorption in Semiconductors}
\label{subsec:pd_maxwell}

The starting point for all optical analysis in this chapter is the macroscopic
form of Maxwell's equations in a non-magnetic, source-free medium. In SI units
and in the frequency domain, obtained by assuming a time dependence
$e^{-i\omega t}$, these read:
\begin{align}
    \nabla \times \tilde{\mathbf{E}}(\mathbf{r},\omega)
        &= i\omega\mu_0\,\tilde{\mathbf{H}}(\mathbf{r},\omega),
        \label{eq:maxwell_curl_e}\\
    \nabla \times \tilde{\mathbf{H}}(\mathbf{r},\omega)
        &= -i\omega\tilde{\varepsilon}(\omega)\,\tilde{\mathbf{E}}(\mathbf{r},\omega),
        \label{eq:maxwell_curl_h}
\end{align}
where $\mu_0$ is the permeability of free space and $\tilde{\varepsilon}(\omega)$
is the complex permittivity. Taking the curl of \cref{eq:maxwell_curl_e} and
substituting \cref{eq:maxwell_curl_h} yields the Helmholtz wave equation:
\begin{equation}
    \nabla^2\tilde{\mathbf{E}}(\mathbf{r},\omega)
    + \tilde{k}^2(\omega)\,\tilde{\mathbf{E}}(\mathbf{r},\omega) = \mathbf{0},
    \label{eq:helmholtz}
\end{equation}
where $\tilde{k}(\omega) = \omega\sqrt{\mu_0\tilde{\varepsilon}(\omega)}$.
Defining the \emph{complex refractive index}:
\begin{equation}
    \tilde{n}(\omega) \equiv \frac{c\,\tilde{k}(\omega)}{\omega}
    = n(\omega) + i\kappa(\omega),
    \label{eq:complex_n}
\end{equation}
the plane-wave solution propagating along $+\hat{z}$ is:
\begin{equation}
    \tilde{\mathbf{E}}(z,\omega)
    = \mathbf{E}_0
      \exp\!\Bigl[i\,n(\omega)\frac{\omega z}{c}\Bigr]
      \exp\!\Bigl[-\kappa(\omega)\frac{\omega z}{c}\Bigr].
    \label{eq:plane_wave}
\end{equation}
The time-averaged Poynting intensity $I(z)\propto|\tilde{\mathbf{E}}|^2$ decays
as $\exp[-2\kappa\omega z/c]$, giving the Beer-Lambert absorption coefficient:
\begin{equation}
    \alpha(\omega) = \frac{2\omega\kappa(\omega)}{c}
    = \frac{4\pi\kappa(\lambda)}{\lambda},
    \label{eq:alpha_kappa}
\end{equation}
so that the Beer-Lambert law:
\begin{equation}
    I(z) = I_0\exp\!\bigl[-\alpha(\lambda)\,z\bigr]
    \label{eq:beer_lambert}
\end{equation}
emerges as a corollary of \cref{eq:helmholtz} rather than an independent
empirical postulate. The absorbed power density follows directly:
\begin{equation}
    \mathcal{P}_{\mathrm{abs}}(\mathbf{r})
    = \alpha(\lambda)\,I(\mathbf{r}).
    \label{eq:abs_density}
\end{equation}

\paragraph{Poynting-vector formulation of the generation rate.}
The optical generation rate $G_{\mathrm{opt}}(\mathbf{r})$ is more generally
obtained from the time-averaged divergence of the Poynting vector. The divergence
theorem applied to the time-averaged Poynting flux
$\langle\mathbf{S}\rangle = \tfrac{1}{2}\mathrm{Re}[\tilde{\mathbf{E}}\times\tilde{\mathbf{H}}^{*}]$
over a closed surface $\partial\Omega$ enclosing a volume $\Omega$ of absorbing
material gives:
\begin{equation}
    \int_\Omega G_{\mathrm{opt}}(\mathbf{r})\,\hbar\omega\,\mathrm{d}V
    = -\oint_{\partial\Omega}
      \langle\mathbf{S}\rangle\cdot\hat{n}\,\mathrm{d}A,
    \label{eq:gopt_poynting}
\end{equation}
so that the local volumetric generation rate is the Poynting-flux sink:
\begin{equation}
    G_{\mathrm{opt}}(\mathbf{r})
    = -\frac{\nabla\cdot\langle\mathbf{S}(\mathbf{r})\rangle}{\hbar\omega}
    = \frac{\omega\varepsilon_0\,\mathrm{Im}[\tilde{\varepsilon}(\mathbf{r})]
           |\tilde{\mathbf{E}}(\mathbf{r})|^2}{2\,\hbar\omega}.
    \label{eq:gopt_poynting_local}
\end{equation}
This reduces identically to \cref{eq:gen_rate} under the plane-wave
approximation, but is the operationally correct definition when the
three-dimensional field contains interference terms arising from modal coupling,
metal reflections, or end-facet diffraction, all of which are present in the
adiabatic taper region of the device studied here.
The FDTD solver employed in \cref{subsec:pd_fdtd} evaluates
$G_{\mathrm{opt}}$ from \cref{eq:gopt_poynting_local}; the scalar Beer-Lambert
law of \cref{eq:beer_lambert} serves only as the analytical validation baseline.

\paragraph{Square-law photodetection and the domain of linear operation.}
For a quasi-monochromatic field at carrier frequency $\omega_0$ with slowly
varying complex envelope $\tilde{s}(t)$:
\begin{equation}
    \tilde{\mathbf{E}}(t)
    = \tfrac{1}{2}\tilde{s}(t)\,e^{-i\omega_0 t}\hat{\mathbf{e}} + \mathrm{c.c.},
    \label{eq:envelope_field}
\end{equation}
the time-averaged Poynting intensity at the detector facet is
$I(t)=|\tilde{s}(t)|^2/(2\eta_0)$, where $\eta_0=\sqrt{\mu_0/\varepsilon_0}$.
The photocurrent produced by the device is therefore:
\begin{equation}
    i_{\mathrm{ph}}(t)
    = \mathcal{R}\cdot P(t)
    = \mathcal{R}\cdot\frac{|\tilde{s}(t)|^2}{2\eta_0},
    \label{eq:square_law}
\end{equation}
which is \emph{linear in optical power} $P(t)$ but \emph{quadratic in optical
field amplitude} $|\tilde{s}(t)|$. When the field contains both a
data-bearing modulated component $\tilde{u}(t)$ and a residual DC carrier $c$,
$\tilde{s}(t)=\tilde{u}(t)+c$, the photocurrent expands as:
\begin{equation}
    i(t) = \mathcal{R}\bigl[
        \underbrace{c^2}_{\text{DC bias}}
       +\underbrace{2c\,\mathrm{Re}\{\tilde{u}(t)\}}_{\text{linear signal term}}
       +\underbrace{|\tilde{u}(t)|^2}_{\text{SSBI}}
    \bigr],
    \label{eq:square_law_expand}
\end{equation}
where SSBI denotes signal-signal beating interference. The linear approximation
$i_{\mathrm{ph}}\approx\mathcal{R}\cdot P_{\mathrm{opt}}$ is valid when the
carrier-to-signal power ratio
$\mathrm{CSPR}\equiv c^2/\langle|\tilde{u}|^2\rangle\gg1$,
so that SSBI is negligible relative to the cross-beat term.

Integrating \cref{eq:abs_density} through an absorber of length $L$ gives the
absorbed fraction:
\begin{equation}
    \mathcal{A}(L) = 1 - \exp\!\bigl[-\alpha(\lambda)\,L\bigr].
    \label{eq:abs_fraction}
\end{equation}

\begin{figure}[H]
    \centering
    \missingfigure[figwidth=0.75\linewidth]{
        Spectral dependence of $\kappa(\lambda)$ and
        $\alpha(\lambda)=4\pi\kappa/\lambda$ in bulk germanium across the
        O- and C-bands. The direct-gap absorption edge near \SI{1550}{\nano\metre}
        and the large $\alpha\approx\SI{5e3}{\per\centi\metre}$ at
        \SI{1310}{\nano\metre} (shaded marker) should be visible. Data source:
        Palik (1985).}
    \caption{Spectral absorption coefficient $\alpha(\lambda)$ and extinction
    coefficient $\kappa(\lambda)$ of bulk germanium, derived via
    \cref{eq:alpha_kappa} from tabulated optical constants. The large $\alpha$
    at \SI{1310}{\nano\metre} reflects direct-gap absorption at the
    $\Gamma$-point discussed in \cref{sec:pd_material}.}
    \label{fig:alpha_kappa}
\end{figure}

% ============================================================
\subsection[Photodetector Figures of Merit]{Photodetector Figures of Merit}
\label{subsec:pd_fom}

Photodetector performance is characterised by a small set of figures of merit
that translate the underlying physics of absorption, transport, and
recombination into engineering-relevant quantities. This section defines these
metrics rigorously and establishes the inter-dependence that governs the
principal design trade-offs encountered in a high-speed Ge-on-Si device
targeting the O-band.

\subsubsection[Quantum Efficiency and Responsivity]{Quantum Efficiency and Responsivity}
\label{subsubsec:pd_qe}

The \emph{external quantum efficiency} $\eta_{\mathrm{ext}}$ is the most
fundamental measure of a photodetector's optical-to-electrical conversion
fidelity. It is defined as the ratio of the number of electron-hole pairs
collected at the device terminals per unit time to the number of photons
incident on the detector per unit time:
\begin{equation}
    \eta_{\mathrm{ext}}
    = \frac{I_{\mathrm{ph}}/q}{P_{\mathrm{opt}}/(\hbar\omega)}
    = \frac{\hbar\omega}{q}\,\mathcal{R},
    \label{eq:eta_ext}
\end{equation}
where $I_{\mathrm{ph}}$ is the photocurrent, $q$ the elementary charge,
$P_{\mathrm{opt}}$ the incident optical power, and $\mathcal{R}$ the
\emph{responsivity} expressed in amperes per watt. The responsivity is a
particularly convenient engineering figure because it is directly measurable
from the slope of a photocurrent-versus-power curve in the linear regime
established by \cref{eq:square_law}:
\begin{equation}
    \mathcal{R} = \frac{I_{\mathrm{ph}}}{P_{\mathrm{opt}}}
    = \frac{\eta_{\mathrm{ext}}\,q}{\hbar\omega}
    = \frac{\eta_{\mathrm{ext}}\,\lambda\,q}{hc}
    \quad [\si{\ampere\per\watt}].
    \label{eq:responsivity_def}
\end{equation}
At the O-band centre wavelength of \SI{1310}{\nano\metre}, the photon energy
is $\hbar\omega = hc/\lambda \approx \SI{0.947}{\electronvolt}$, and the
theoretical upper bound on responsivity, achieved when $\eta_{\mathrm{ext}}=1$,
is:
\begin{equation}
    \mathcal{R}_{\max}
    = \frac{q\lambda}{hc}
    \approx \SI{1.057}{\ampere\per\watt} \quad (\lambda = \SI{1310}{\nano\metre}).
    \label{eq:R_max}
\end{equation}
The external quantum efficiency can be factored into contributions from
coupling, absorption, and collection:
\begin{equation}
    \eta_{\mathrm{ext}}
    = \eta_{\mathrm{coup}}\,\eta_{\mathrm{abs}}\,\eta_{\mathrm{coll}},
    \label{eq:eta_product}
\end{equation}
where $\eta_{\mathrm{coup}}$ accounts for optical coupling losses at the
silicon-waveguide-to-germanium interface (mode mismatch, reflection),
$\eta_{\mathrm{abs}} = \mathcal{A}(L) = 1-e^{-\alpha L}$ is the absorbed
fraction derived in \cref{eq:abs_fraction}, and $\eta_{\mathrm{coll}}$
captures the fraction of photogenerated carriers that reach the contacts
without recombining. In a well-depleted PIN structure operated at sufficient
reverse bias, $\eta_{\mathrm{coll}}\rightarrow 1$ because field-driven drift
removes carriers on a timescale $\tau_{\mathrm{tr}} \ll \tau_{\mathrm{SRH}}$.
The product $\eta_{\mathrm{coup}}\,\eta_{\mathrm{abs}}$ is therefore the
dominant lever in device design: the former is addressed through taper geometry
and modal overlap engineering, while the latter dictates the minimum absorber
length. The auxiliary MATLAB model used in this repository takes
$\alpha\approx\SI{7e3}{\per\centi\metre}$ at
\SI{1310}{\nano\metre}, corresponding to a bulk absorption length of
$\alpha^{-1}\approx\SI{1.4}{\micro\metre}$ before waveguide-confinement and
mode-overlap factors are applied. The final absorber length is therefore not
selected from the bulk Beer-Lambert estimate alone; it is selected from the
FDTD taper-coupled length sweep.

\paragraph{Internal quantum efficiency and the collection efficiency limit.}
It is instructive to distinguish the \emph{internal} quantum efficiency
$\eta_{\mathrm{int}}$, which excludes coupling and reflection losses and
characterises solely the conversion of absorbed photons to collected
charge~\cite{SzeNg2006}. In the limit where all photogeneration is confined to
the fully depleted intrinsic region of length $W_i$ and drift transport is
complete, $\eta_{\mathrm{int}} = 1$. Departures from this ideal arise when
photogeneration extends into quasi-neutral regions outside the depletion zone,
where diffusion is the dominant transport mechanism. The diffusion minority
carrier lifetime $\tau_n$ in the $p$-doped anode sets a characteristic
diffusion length $L_n = \sqrt{D_n\tau_n}$; only carriers generated within $L_n$
of the depletion boundary contribute to the photocurrent. Carriers generated
beyond this distance recombine before reaching the junction. This mechanism is
especially pertinent in the O-band context, where the large absorption depth
at wavelengths approaching the indirect-gap edge can push a fraction of the
generation profile into the quasi-neutral region if the intrinsic layer is not
sufficiently thick or the reverse bias is insufficient to extend the depletion
width. The trade-off is explicit: increasing $W_i$ increases $\eta_{\mathrm{int}}$
but simultaneously increases the transit time $\tau_{\mathrm{tr}} = W_i / v_{\mathrm{d,sat}}$
and the junction capacitance $C_j = \varepsilon_{\mathrm{Ge}}\,A_j / W_i$, both
of which degrade bandwidth as formalised in \cref{subsec:pd_bw}.

\subsubsection[Bandwidth and Speed Limitations]{Bandwidth and Speed Limitations}
\label{subsubsec:pd_bw}

The small-signal \SI{3}{dB} bandwidth $f_{3\mathrm{dB}}$ of a photodetector is
the electrical frequency at which the detector's photocurrent response is
attenuated by \SI{3}{dB} relative to its low-frequency value. Two fundamentally
distinct physical mechanisms limit this bandwidth and must be engineered
simultaneously for high-speed operation.

\paragraph{Transit-time bandwidth.}
When an electron-hole pair is generated inside the intrinsic region and swept
towards the contacts by the applied field, the terminal current induced via the
Ramo-Shockley theorem is non-zero for the duration of transit. A pair
generated at depth $z$ contributes a rectangular current pulse of duration
$\tau_e(z) = z/v_{\mathrm{d,sat},e}$ for the electron and
$\tau_h(z) = (W_i - z)/v_{\mathrm{d,sat},h}$ for the hole, as expressed in
\cref{eq:tau_depth,eq:hz_local}. The frequency response of a single pair
follows from the Fourier transform of this pulse pair, \cref{eq:hz_freq},
and upon averaging uniformly over all generation depths within $[0,W_i]$
via \cref{eq:htt_integral}, the transit-time transfer function takes the form
of a weighted sinc envelope. In the practically relevant limit where the hole
branch---slower due to $v_{\mathrm{d,sat},h} < v_{\mathrm{d,sat},e}$ in
germanium---dominates the roll-off, the \SI{3}{dB} transit-time bandwidth is
well approximated by:
\begin{equation}
    f_{\mathrm{tt}}
    \approx \frac{0.443\,v_{\mathrm{d,sat},h}}{W_i},
    \label{eq:bw_tt}
\end{equation}
where $v_{\mathrm{d,sat},h} \approx \SI{6.0e6}{\centi\metre\per\second}$ for
holes in germanium at fields exceeding approximately
\SI{10}{\kilo\volt\per\centi\metre}~\cite{SzeNg2006}. The numerical prefactor
0.443 arises from the sinc-function roll-off $|\mathrm{sinc}(\omega\tau/2)|=1/\sqrt{2}$
at $\omega\tau_{\mathrm{tr}} = 2.78$. \Cref{eq:bw_tt} immediately establishes
that transit-time bandwidth is inversely proportional to the intrinsic-layer
thickness: halving $W_i$ doubles $f_{\mathrm{tt}}$.

\paragraph{RC-limited bandwidth.}
The second mechanism is purely electrical in origin. The detector junction
presents a capacitance $C_j$ to the external circuit, whose total load
resistance is $R_{\mathrm{tot}} = R_s + R_L$, where $R_s$ is the device
series resistance (contact and spreading resistance) and $R_L = \SI{50}{\ohm}$
is the standard termination impedance in high-speed link applications. The
resulting single-pole low-pass response, \cref{eq:hrc_transfer}, gives the
RC-limited bandwidth:
\begin{equation}
    f_{\mathrm{RC}}
    = \frac{1}{2\pi\,R_{\mathrm{tot}}\,C_j}.
    \label{eq:bw_rc}
\end{equation}
For a planar junction of area $A_j = W_j \times L$ and intrinsic thickness
$W_i$, the depletion capacitance is:
\begin{equation}
    C_j = \frac{\varepsilon_0\,\varepsilon_{\mathrm{Ge}}\,A_j}{W_i},
    \label{eq:cj_planar}
\end{equation}
where $\varepsilon_{\mathrm{Ge}} \approx 16.0$ is the static relative
permittivity of germanium. Substituting into \cref{eq:bw_rc} makes the
competing roles of $W_i$ explicit: increasing $W_i$ reduces $C_j$ and raises
$f_{\mathrm{RC}}$, but simultaneously reduces $f_{\mathrm{tt}}$ via
\cref{eq:bw_tt}. This \emph{transit-time/RC trade-off} is the central design
tension in the intrinsic-layer thickness optimisation.

\paragraph{Combined bandwidth.}
Because the transit-time and RC contributions to bandwidth act as independent
roll-off mechanisms in cascade, the total \SI{3}{dB} bandwidth of the device
is:
\begin{equation}
    \frac{1}{f_{3\mathrm{dB}}^2}
    \approx \frac{1}{f_{\mathrm{tt}}^2} + \frac{1}{f_{\mathrm{RC}}^2},
    \label{eq:bw_combined}
\end{equation}
where the quadrature sum is an approximation that is exact when both
sub-responses are first-order Butterworth poles and remains accurate to within
a few percent for the sinc-envelope transit-time case~\cite{Beling2014}.
The full factorised transfer function is given by \cref{eq:hpd_factorised}.
Optimising $W_i$ to maximise $f_{3\mathrm{dB}}$ subject to a fixed device
footprint then reduces to finding the value of $W_i$ satisfying
$\partial f_{3\mathrm{dB}} / \partial W_i = 0$, which for $R_{\mathrm{tot}} =
\SI{50}{\ohm}$ typically places the optimum in the range
$W_i\in[\SI{0.5}{\micro\metre},\,\SI{1.5}{\micro\metre}]$ for a
waveguide-integrated Ge-on-Si detector with a junction width of order
\SI{1}{\micro\metre} targeting bandwidths in the range \SIrange{50}{70}{\giga\hertz}.

\begin{figure}[H]
    \centering
    \missingfigure[figwidth=0.80\linewidth]{
        Combined bandwidth as a function of intrinsic-layer thickness $W_i$
        from \cref{eq:bw_combined}, showing the transit-time limit
        $f_{\mathrm{tt}}(W_i)$ and RC limit $f_{\mathrm{RC}}(W_i)$ as
        individual curves and their quadrature combination. The optimal
        $W_i^*$ is identified at the peak of $f_{3\mathrm{dB}}$.}
    \caption{Transit-time and RC bandwidth components as a function of the
    intrinsic-layer thickness $W_i$ for a Ge-on-Si waveguide detector with
    $A_j = \SI{1}{\micro\metre}\times\SI{15}{\micro\metre}$ and
    $R_{\mathrm{tot}} = \SI{50}{\ohm}$. The optimal thickness $W_i^*$ yields
    the maximum combined bandwidth, illustrating the fundamental transit-time/RC
    trade-off.}
    \label{fig:bw_vs_wi}
\end{figure}

\subsubsection[Noise Mechanisms and Signal-to-Noise Ratio]{Noise Mechanisms and Signal-to-Noise Ratio}
\label{subsubsec:pd_noise}

Noise ultimately determines the minimum detectable optical power and, in
conjunction with bandwidth, the achievable bit-error rate (BER) in a digital
link. Three noise sources are relevant in a reverse-biased PIN
photodetector~\cite{SzeNg2006}: shot noise from both signal and dark
photocurrents, and thermal noise from the load resistance.

\paragraph{Shot noise.}
Shot noise originates from the quantum-mechanical discreteness of the
photodetection process: each photon-induced electron-hole pair is generated at a
statistically independent random instant, following Poisson statistics. The
two-sided power spectral density (PSD) of the resulting current fluctuation
$\delta i_{\mathrm{sh}}$ is flat (white) up to the detector bandwidth and is
given by:
\begin{equation}
    S_{\mathrm{sh}}(f) = 2q\,(I_{\mathrm{ph}} + I_d),
    \label{eq:shot_psd}
\end{equation}
where $I_{\mathrm{ph}} = \mathcal{R}\,P_{\mathrm{opt}}$ is the signal
photocurrent and $I_d$ is the dark current. The mean-square shot-noise current
within a measurement bandwidth $B$ is therefore:
\begin{equation}
    \langle i_{\mathrm{sh}}^2 \rangle = 2q\,(I_{\mathrm{ph}} + I_d)\,B.
    \label{eq:shot_noise_var}
\end{equation}
In the context of a Ge-on-Si detector, the dark current $I_d$ is substantially
larger than in silicon counterparts because of the high density of SRH trap
centres introduced by threading dislocations during heteroepitaxial growth. A
representative dark current of order \SIrange{10}{100}{\nano\ampere} at
$V_{\mathrm{bias}} = \SI{-1}{\volt}$ for a device of area
$A_j\approx\SI{15}{\micro\metre\squared}$ implies a dark-current shot-noise PSD
that can rival the signal shot noise at low received powers. Minimising $I_d$
through lifetime engineering and careful epitaxial quality control is therefore
not merely a static power concern but directly impacts the receiver noise floor.

\paragraph{Thermal (Johnson) noise.}
The load resistor $R_L$ at temperature $T$ introduces a Johnson noise current
with PSD:
\begin{equation}
    S_{\mathrm{th}}(f) = \frac{4k_B T}{R_L},
    \label{eq:thermal_psd}
\end{equation}
which is also white within the bandwidth of interest. The mean-square thermal
noise current in bandwidth $B$ is:
\begin{equation}
    \langle i_{\mathrm{th}}^2 \rangle = \frac{4k_B T}{R_L}\,B.
    \label{eq:thermal_noise_var}
\end{equation}
At room temperature ($T = \SI{300}{\kelvin}$) and standard termination
$R_L = \SI{50}{\ohm}$, the thermal noise current density is
$\sqrt{S_{\mathrm{th}}} \approx \SI{18}{\pico\ampere\per\hertz^{1/2}}$, which
is typically the dominant noise source at low optical power levels in a
direct-detection receiver without optical preamplification.

\paragraph{Signal-to-noise ratio and noise-equivalent power.}
The electrical signal-to-noise ratio (SNR) for a direct-detection receiver
combining shot noise and thermal noise is:
\begin{equation}
    \mathrm{SNR}
    = \frac{I_{\mathrm{ph}}^2}
           {\langle i_{\mathrm{sh}}^2 \rangle + \langle i_{\mathrm{th}}^2 \rangle}
    = \frac{(\mathcal{R}\,P_{\mathrm{opt}})^2}
           {2q(\mathcal{R}\,P_{\mathrm{opt}} + I_d)B + 4k_BT B/R_L}.
    \label{eq:snr}
\end{equation}
The \emph{noise-equivalent power} (NEP) is defined as the optical input power
at which SNR $= 1$ in a bandwidth $B = \SI{1}{\hertz}$, providing a bandwidth-
independent figure of merit for detector sensitivity:
\begin{equation}
    \mathrm{NEP} = \frac{\sqrt{2q\,I_d + 4k_BT/R_L}}{\mathcal{R}}
    \quad [\si{\watt\per\hertz^{1/2}}].
    \label{eq:nep}
\end{equation}
\Cref{eq:nep} explicitly couples the dark-current magnitude, the responsivity,
and the load impedance into a single sensitivity figure. For the Ge-on-Si
device characterised in this work, the dark-current term dominates at room
temperature and typical bias conditions, so that reducing $I_d$ and increasing
$\mathcal{R}$ are the two principal levers for sensitivity improvement. In the
shot-noise-limited regime, where $I_d\rightarrow0$ and $R_L\rightarrow\infty$
(achieved in practice with a transimpedance amplifier of very large
input impedance), the NEP reduces to the quantum-limited floor:
\begin{equation}
    \mathrm{NEP}_{\mathrm{shot}} = \frac{\sqrt{2q\,I_{\mathrm{ph}}}}{\mathcal{R}}
    \Big|_{I_{\mathrm{ph}}\,\rightarrow\,0}
    \xrightarrow{\eta_{\mathrm{ext}}\rightarrow 1}
    \sqrt{2\hbar\omega}\,.
    \label{eq:nep_shot_limit}
\end{equation}
For the O-band, this evaluates to $\mathrm{NEP}_{\mathrm{shot}} \approx
\SI{5.5e-19}{\watt\per\hertz^{1/2}}$, a value that practical devices approach
asymptotically as dark current and excess noise are suppressed. The derivation
of the dark-current model $I_d(V_{\mathrm{bias}})$ used in \cref{eq:nep} is
presented in \cref{subsec:pd_charge}, where the surface and bulk SRH components
are separated and independently calibrated.

% ============================================================
\subsection[Carrier Photogeneration and Transport]{Carrier Photogeneration and Transport}
\label{subsec:pd_generation}

The absorbed photon creates an electron-hole pair. Whether that pair is swept
to the contacts and contributes to the photocurrent, or instead recombines
within the absorber and is lost, depends on the interplay of the electric
field, the carrier mobilities, and the defect density of the semiconductor.

The volumetric optical generation rate $G_{\mathrm{opt}}(\mathbf{r})$ quantifies
the number of electron-hole pairs created per unit volume per unit time:
\begin{equation}
    G_{\mathrm{opt}}(\mathbf{r})
    = \frac{\mathcal{P}_{\mathrm{abs}}(\mathbf{r})}{\hbar\omega}
    = \frac{\alpha(\lambda)\,I(\mathbf{r})}{\hbar\omega}
    \quad [\si{\per\centi\metre\cubed\per\second}].
    \label{eq:gen_rate}
\end{equation}
Because $I(\mathbf{r})$ decays exponentially along the propagation direction via
\cref{eq:beer_lambert}, $G_{\mathrm{opt}}$ is spatially non-uniform: it is
concentrated near the input facet of the absorber and falls off with an
effective length set by the Ge material absorption, the optical confinement
factor, and the simulated taper-to-Ge mode overlap.

\paragraph{Drift--diffusion transport.}
The self-consistent carrier distribution and electrostatic field are governed
by the coupled Poisson and drift-diffusion equations~\cite{SzeNg2006,Yariv1997}.
Poisson's equation is:
\begin{equation}
    \nabla\cdot\!\left[\varepsilon(\mathbf{r})\,\nabla\phi(\mathbf{r},t)\right]
    = -q\!\left[p - n + N_D^+ - N_A^-\right],
    \label{eq:poisson}
\end{equation}
and the electron and hole current densities in the drift-diffusion approximation
are:
\begin{align}
    \mathbf{J}_n &= q\mu_n\,n\,\mathbf{E} + qD_n\,\nabla n,
        \label{eq:Jn}\\
    \mathbf{J}_p &= q\mu_p\,p\,\mathbf{E} - qD_p\,\nabla p,
        \label{eq:Jp}
\end{align}
where the carrier diffusivities follow the Einstein relation~\cite{Hamadou2023}:
\begin{equation}
    D_{n,p} = \mu_{n,p}\,\frac{k_B T}{q}.
    \label{eq:einstein}
\end{equation}
The general, time-dependent carrier continuity equations are:
\begin{align}
    \frac{\partial n}{\partial t}
        &= \frac{1}{q}\nabla\cdot\mathbf{J}_n
           - U_{\mathrm{net}} + G_{\mathrm{opt}},
        \label{eq:cont_n_transient}\\
    \frac{\partial p}{\partial t}
        &= -\frac{1}{q}\nabla\cdot\mathbf{J}_p
           - U_{\mathrm{net}} + G_{\mathrm{opt}},
        \label{eq:cont_p_transient}
\end{align}
where $U_{\mathrm{net}}=U_{\mathrm{SRH}}+U_{\mathrm{Auger}}+U_{\mathrm{rad}}$.
Under \emph{steady-state} CW illumination, $\partial/\partial t\equiv0$:
\begin{align}
    \frac{1}{q}\nabla\cdot\mathbf{J}_n &= U_{\mathrm{net}} - G_{\mathrm{opt}},
        \label{eq:cont_n}\\
    \frac{1}{q}\nabla\cdot\mathbf{J}_p &= -(U_{\mathrm{net}} - G_{\mathrm{opt}}).
        \label{eq:cont_p}
\end{align}
The complete self-consistent system, from \cref{eq:poisson} through
\cref{eq:cont_p}, is the foundational PDE set solved numerically by the
commercial finite-element CHARGE solver and reduced in MATLAB for the
geometry-derived bandwidth and noise models.

The central design consequence of the drift-diffusion picture is the distinction
between drift-dominated and diffusion-dominated carrier motion.
Drift transport, which is field-driven and fast, requires a sufficiently large
reverse-bias field to sweep carriers across the intrinsic region before they
can recombine. Diffusion transport, which is concentration-gradient-driven and
inherently slower, dominates in the quasi-neutral regions outside the depletion
zone and is the primary cause of slow response tails in devices with incomplete
or absent depletion. A reverse-biased PIN junction suppresses diffusion by
confining photogeneration to the depleted region, making drift the dominant
transport mechanism and the carrier transit time the key speed parameter.

\paragraph{Velocity saturation and field-dependent mobility.}
An important refinement to the linear drift picture arises at the field
magnitudes present in a reverse-biased Ge PIN junction. The carrier drift
velocity does not increase indefinitely with electric field but saturates at
$v_{\mathrm{d,sat}}$ when the field exceeds a critical value
$\mathcal{E}_c$. A widely adopted empirical model for the field-dependent
drift velocity is the Caughey-Thomas expression~\cite{SzeNg2006}:
\begin{equation}
    v_{n,p}(\mathcal{E})
    = \frac{\mu_{n,p}^{(0)}\,\mathcal{E}}
           {\left[1 + \left(\mu_{n,p}^{(0)}\mathcal{E}/v_{\mathrm{d,sat},n/p}
           \right)^{\beta_{n,p}}\right]^{1/\beta_{n,p}}},
    \label{eq:caughey_thomas}
\end{equation}
where $\mu_{n,p}^{(0)}$ is the low-field mobility, $v_{\mathrm{d,sat},n/p}$
is the saturation velocity, and $\beta_{n,p}$ is an empirical shape parameter.
For germanium at \SI{300}{\kelvin}, representative values are
$\mu_n^{(0)} \approx \SI{3900}{\centi\metre\squared\per\volt\per\second}$,
$\mu_p^{(0)} \approx \SI{1900}{\centi\metre\squared\per\volt\per\second}$,
$v_{\mathrm{d,sat},e} \approx \SI{6.0e6}{\centi\metre\per\second}$, and
$v_{\mathrm{d,sat},h} \approx \SI{6.0e6}{\centi\metre\per\second}$, with
$\beta\approx 2$ for electrons and $\beta\approx 1$ for holes~\cite{SzeNg2006}.
At a reverse bias of \SI{2}{\volt} across a \SI{1}{\micro\metre} intrinsic
layer, the mean field is \SI{20}{\kilo\volt\per\centi\metre}, which
substantially exceeds the critical field
$\mathcal{E}_c = v_{\mathrm{d,sat}}/\mu^{(0)}
\approx \SIrange{1}{3}{\kilo\volt\per\centi\metre}$ and places the device
firmly in the velocity-saturated regime. The constant velocity assumption
adopted in \cref{eq:bw_tt} and the analytical transfer function of
\cref{eq:jt_analytical} is therefore well justified at standard operating
bias, and the transit time is accurately estimated as
$\tau_{\mathrm{tr}} = W_i/v_{\mathrm{d,sat}}$ rather than the low-field
value $W_i^2/(\mu V)$.

\paragraph{Analytical photocurrent transfer function.}
In the frequency domain, neglecting diffusion and assuming a uniform electric
field at saturation value throughout $z\in[0,W_i]$, the Fourier amplitude of
the terminal current density admits a closed-form expression:
\begin{align}
    J_t(\omega)
    &= \frac{q G_0}{\alpha W_i}\frac{hc}{q\lambda}
       \eta(1-R_{\mathrm{refl}})P_{\mathrm{opt}}(\omega)\nonumber\\
    &\quad\times\!
      \left\{
        \frac{e^{-\alpha W_i}-1}{\alpha W_i(\alpha W_i - j\omega\tau_{\mathrm{tr,h}})}
        + e^{-\alpha W_i}
          \frac{e^{j\omega\tau_{\mathrm{tr,h}}}-1}
               {j\omega\tau_{\mathrm{tr,h}}(\alpha W_i - j\omega\tau_{\mathrm{tr,h}})}
      \right.\nonumber\\
    &\quad\quad\left.
        + \frac{1 - e^{-\alpha W_i}}{\alpha W_i(j\omega\tau_{\mathrm{tr,e}} + \alpha W_i)}
        + \frac{1 - e^{j\omega\tau_{\mathrm{tr,e}}}}{j\omega\tau_{\mathrm{tr,e}}
          (j\omega\tau_{\mathrm{tr,e}} + \alpha W_i)}
      \right\},
    \label{eq:jt_analytical}
\end{align}
where $\tau_{\mathrm{tr,e}}=W_i/v_{\mathrm{d,sat},e}$ and
$\tau_{\mathrm{tr,h}}=W_i/v_{\mathrm{d,sat},h}$ are the electron and hole
transit times. In the limits $\omega\tau_{\mathrm{tr,h}}\ll1$ and
$\alpha W_i\gg1$, \cref{eq:jt_analytical} reduces to the rectangular-pulse
sinc-function roll-off yielding the bandwidth formula of \cref{eq:bw_tt}.

\begin{figure}[H]
    \centering
    \missingfigure[figwidth=0.86\linewidth]{
        Two-part transport figure. Left: a carrier pair generated inside the
        intrinsic layer with electron and hole transit paths to the contacts.
        Right: equivalent-circuit derivation flow
        $h_z(t)\rightarrow H_{\mathrm{tt}}(\omega)\rightarrow
        H_{\mathrm{RC}}(\omega)\rightarrow H_{\mathrm{PD}}(\omega)$.}
    \caption{Microscopic-to-circuit derivation path for the small-signal transfer
    function. The transit-time response originates from the finite carrier flight
    time inside the intrinsic Ge region, while the RC response originates from
    the external electrical loading seen by the junction.}
    \label{fig:transfer_derivation}
\end{figure}

\paragraph{Derivation of the factorised small-signal transfer function.}
Consider an electron-hole pair generated at depth $z\in[0,W_i]$ inside the
intrinsic layer. Under the uniform-field, velocity-saturated approximation, the
electron and hole transit times are:
\begin{equation}
    \tau_e(z)=\frac{z}{v_{\mathrm{d,sat},e}},
    \qquad
    \tau_h(z)=\frac{W_i-z}{v_{\mathrm{d,sat},h}}.
    \label{eq:tau_depth}
\end{equation}
By the Ramo-Shockley theorem, the terminal current induced by each carrier is
non-zero only while the carrier is moving, giving the local impulse response:
\begin{equation}
    h_z(t)
    = \frac{qv_{\mathrm{d,sat},e}}{W_i}\,
      u(t)\,u\!\left[\tau_e(z)-t\right]
    + \frac{qv_{\mathrm{d,sat},h}}{W_i}\,
      u(t)\,u\!\left[\tau_h(z)-t\right],
    \label{eq:hz_local}
\end{equation}
where $u(t)$ is the Heaviside function. Fourier transformation yields:
\begin{equation}
    H_z(\omega)
    = \frac{q}{W_i}\left[
        \frac{1-e^{-j\omega\tau_e(z)}}{j\omega}
      + \frac{1-e^{-j\omega\tau_h(z)}}{j\omega}
      \right].
    \label{eq:hz_freq}
\end{equation}
If the optical generation is approximately uniform across the depletion depth,
the transit response follows from averaging \cref{eq:hz_freq} over $z$:
\begin{equation}
    H_{\mathrm{tt}}(\omega)\propto
    \int_0^{W_i} H_z(\omega)\,\frac{\mathrm{d}z}{W_i}.
    \label{eq:htt_integral}
\end{equation}
The integral evaluates to a weighted sum of electron and hole sinc terms. In
the practically relevant limit where the slower hole branch dominates the
\SI{3}{dB} roll-off, the envelope reduces to the single-sinc form used in
\cref{eq:H_pd}. The electrical loading is then appended through the junction
capacitance and total series resistance. For a current source driving
$C_j$ into $R_{\mathrm{tot}}=R_s+R_L$:
\begin{equation}
    H_{\mathrm{RC}}(\omega)
    = \frac{1}{1+j\omega R_{\mathrm{tot}}C_j}
    = \frac{1}{1+j\omega/\omega_{\mathrm{RC}}},
    \label{eq:hrc_transfer}
\end{equation}
with $\omega_{\mathrm{RC}}=(R_{\mathrm{tot}}C_j)^{-1}$. The total response is
the product:
\begin{equation}
    H_{\mathrm{PD}}(\omega)
    = H_{\mathrm{tt}}(\omega)\,H_{\mathrm{RC}}(\omega),
    \label{eq:hpd_factorised}
\end{equation}
which is the physical content hidden inside the compact form of
\cref{eq:H_pd}: transport determines the sinc-like envelope, while the external
electrical boundary condition determines the low-pass pole.

\paragraph{Scharfetter-Gummel discretisation and numerical stability.}
The finite-element discretisation of \cref{eq:Jn,eq:Jp} on a non-uniform mesh
faces a well-known numerical instability when the local P\'eclet number
$\mathrm{Pe}=q|\mathbf{E}|\Delta x/(2kT)$ exceeds unity. The
Scharfetter-Gummel scheme approximates the current between adjacent mesh nodes
$i$ and $i+1$ as:
\begin{equation}
    J_{n,i+1/2}
    = \frac{qD_n}{\Delta x}\!\left[
        B\!\left(\frac{q\Delta\phi}{kT}\right)n_{i+1}
      - B\!\left(-\frac{q\Delta\phi}{kT}\right)n_i
      \right],
    \label{eq:sg_scheme}
\end{equation}
where $B(u)=u/(e^u-1)$ is the Bernoulli function and $\Delta\phi=\phi_{i+1}-\phi_i$.
This scheme recovers centred differences for $|u|\ll1$ (low-field) and
upstream weighting for $|u|\gg1$ (high-field), providing second-order accuracy
and unconditional positivity across the \SI{30}{\kilo\volt\per\centi\metre}
field regime of the present device.

% ============================================================
\subsection[The PIN Junction: Electrostatics and Operating Regime]%
           {The PIN Junction: Electrostatics and Operating Regime}
\label{subsec:pd_pin}

The vertical PIN architecture---a lightly doped or nominally intrinsic
germanium absorber sandwiched between heavily doped $p^{++}$ and $n^{++}$ contact
layers---is the canonical structure for high-speed photodetection. Its
central merit is the ability to sustain a large, spatially uniform electric
field across the entire absorber volume under modest reverse bias, simultaneously
achieving full depletion, velocity-saturated carrier transport, and minimal
free-carrier absorption.

\subsubsection[Depletion Approximation and Built-in Potential]{Depletion Approximation and Built-in Potential}
\label{subsubsec:pd_depletion}

In the absence of applied bias, the Fermi levels of the $p^{++}$ and $n^{++}$
regions equilibrate by charge redistribution, establishing a built-in potential:
\begin{equation}
    V_{\mathrm{bi}}
    = \frac{k_B T}{q}\ln\!\left(\frac{N_A\,N_D}{n_i^2}\right),
    \label{eq:vbi}
\end{equation}
where $N_A$ and $N_D$ are the acceptor and donor concentrations on the $p$ and
$n$ sides respectively, and $n_i$ is the intrinsic carrier concentration of
germanium. At \SI{300}{\kelvin}, $n_i^{\mathrm{Ge}} \approx
\SI{2.4e13}{\per\centi\metre\cubed}$, approximately three orders of magnitude
larger than in silicon, which has important consequences for dark current as
discussed in \cref{subsec:pd_recombination}. For typical contact doping levels
of $N_{A,D} \approx \SI{e19}{\per\centi\metre\cubed}$, \cref{eq:vbi} gives
$V_{\mathrm{bi}} \approx \SI{0.51}{\volt}$.

Under the depletion approximation, the space charge in the $p^{++}$ and
$n^{++}$ layers is assumed to be that of fully ionised impurities with sharp
depletion edges. In the ideal case where the intrinsic layer is fully depleted
at zero bias by the built-in field, the charge neutrality condition reads:
\begin{equation}
    N_A\,x_p = N_D\,x_n,
    \label{eq:charge_neutral}
\end{equation}
where $x_p$ and $x_n$ are the partial penetrations of the depletion region
into the $p^{++}$ and $n^{++}$ layers, respectively. Since
$N_{A,D} \gg n_i$, the depletion extension into the heavily doped contact
layers is negligible ($x_p, x_n \ll W_i$), and the electric field profile
is approximately uniform across the intrinsic layer with peak value:
\begin{equation}
    \mathcal{E}_{\max}
    = \frac{q N_A x_p}{\varepsilon_0 \varepsilon_{\mathrm{Ge}}}
    \approx \frac{V_{\mathrm{bi}} + |V_{\mathrm{bias}}|}{W_i},
    \label{eq:E_max}
\end{equation}
where the second equality applies in the fully depleted regime. For
$V_{\mathrm{bias}} = \SI{-2}{\volt}$ and $W_i = \SI{1}{\micro\metre}$,
\cref{eq:E_max} gives $\mathcal{E}_{\max} \approx \SI{25}{\kilo\volt\per\centi\metre}$,
confirming velocity saturation. The flatness of the field profile is a key
advantage of the PIN geometry over a simple $pn$ junction: it ensures that
every photogenerated carrier, regardless of its generation depth within $[0,W_i]$,
experiences the same drift force and therefore the same transit time, making
\cref{eq:tau_depth} well-conditioned for all $z$.

\subsubsection[Dark Current Mechanisms and Decomposition]{Dark Current Mechanisms and Decomposition}
\label{subsubsec:pd_idark}

The dark current $I_d$ under reverse bias is the primary figure of merit for
detector noise and static power dissipation. In a Ge-on-Si PIN device, it
receives contributions from both the junction bulk volume and the mesa
perimeter, and it is important to distinguish these because they scale
differently with device geometry:
\begin{equation}
    I_d = J_{\mathrm{bulk}}\,A_j + J_{\mathrm{surf}}\,P_j,
    \label{eq:idark_decomposed}
\end{equation}
where $J_{\mathrm{bulk}}$ is the bulk dark-current density, $J_{\mathrm{surf}}$
is an effective surface dark-current line density, and $P_j$ is the mesa
perimeter~\cite{SzeNg2006}. The bulk component in the fully depleted intrinsic
region is dominated by SRH thermal generation (as established in
\cref{subsec:pd_recombination}) and is given by:
\begin{equation}
    J_{\mathrm{bulk}}
    = q\,U_{\mathrm{SRH,dep}}\,W_i
    = \frac{q\,n_i\,W_i}{2\tau_0},
    \label{eq:jdark_bulk}
\end{equation}
which scales linearly with $W_i$ and inversely with the effective SRH lifetime
$\tau_0$. The surface component arises from interface states at the Ge/SiO$_2$
sidewall and Ge/Si heterointerface and is written as:
\begin{equation}
    J_{\mathrm{surf}}
    = q\,v_s\,n_i,
    \label{eq:jdark_surf}
\end{equation}
where $v_s$ is the effective surface recombination velocity. Because $n_i$ of
germanium is orders of magnitude larger than that of silicon, both
$J_{\mathrm{bulk}}$ and $J_{\mathrm{surf}}$ in Ge-on-Si are intrinsically
larger than in silicon-only detectors by the same factor, placing a fundamental
material constraint on achievable dark current density. As device dimensions
are scaled to improve bandwidth---reducing $A_j$ and $P_j$ together but with
$A_j/P_j \propto W_j$ decreasing---the surface term becomes proportionally
more important, and surface passivation through dielectric encapsulants such as
Al$_2$O$_3$ or GeO$_x$ becomes critical. The decomposition of \cref{eq:idark_decomposed}
is exploited in the calibration procedure of \cref{subsec:pd_charge}, where
measurements on devices of varying junction width $W_j$ are used to extract
$J_{\mathrm{bulk}}$ and $J_{\mathrm{surf}}$ independently.

% ============================================================
\subsection[Recombination and Loss Mechanisms]{Recombination and Loss Mechanisms}
\label{subsec:pd_recombination}

Three recombination mechanisms are quantitatively relevant in the germanium
absorber~\cite{SzeNg2006,Basu2003}.

\emph{Shockley-Read-Hall (SRH) recombination}, mediated by deep-level trap
states, is~\cite{Hamadou2023}:
\begin{equation}
    U_{\mathrm{SRH}}
    = \frac{np - n_i^2}{\tau_p\!\left(n + n_1\right)
      + \tau_n\!\left(p + p_1\right)},
    \label{eq:srh}
\end{equation}
where the effective trap-state carrier densities are:
\begin{align}
    n_1 &= n_i\exp\!\left[\frac{E_t-E_i}{k_BT}\right], &
    p_1 &= n_i\exp\!\left[\frac{E_i-E_t}{k_BT}\right].
    \label{eq:np1_trap}
\end{align}

\paragraph{Fossum doping-dependent SRH lifetime.}
In the presence of a background doping $N_A+N_D$, the SRH carrier lifetime is
not constant but falls with increasing impurity concentration. The Fossum
model~\cite{Hamadou2023} captures this dependence as:
\begin{equation}
    \tau_{\mathrm{SRH},n/p}
    = \frac{\tau^{(0)}_{\mathrm{SRH},n/p}}
           {a_1 + a_2\,\mathcal{N}_{n/p}
            + a_3\,\mathcal{N}_{n/p}^{\,a_4}},
    \quad
    \mathcal{N}_{n/p} = \frac{N_A+N_D}{N^{(\mathrm{ref})}_{n/p}},
    \label{eq:fossum}
\end{equation}
where $\tau^{(0)}_{\mathrm{SRH},n/p}$ is the low-concentration lifetime and
$a_1,\ldots,a_4$ are fitting parameters. In the limit $a_2=0$,
$a_1=a_3=a_4=1$, \cref{eq:fossum} reduces to the standard
Shockley-Read-Hall result $\tau=\tau^{(0)}/(1+N/N^{(\mathrm{ref})})$.
For the heavily doped $n^{++}$-Ge cathode at
$N_D=\SI{e19}{\per\centi\metre\cubed}$, \cref{eq:fossum} gives a lifetime
reduction of approximately one order of magnitude relative to $\tau^{(0)}$,
confirming that carrier recombination is concentrated in the contact region
rather than the intrinsic absorber.
Within the fully depleted intrinsic region at the midgap trap condition
$E_t=E_i$, \cref{eq:srh} reduces to:
\begin{equation}
    U_{\mathrm{SRH,dep}} \approx \frac{n_i}{\tau_n + \tau_p}
    \equiv \frac{n_i}{2\tau_0},
    \label{eq:srh_dep}
\end{equation}
where $\tau_0=(\tau_n+\tau_p)/2$. SRH is the dominant mechanism in Ge-on-Si
because the threading dislocations formed during heteroepitaxial growth act as
efficient midgap trap centres.

\emph{Auger recombination} is~\cite{Hamadou2023}:
\begin{equation}
    U_{\mathrm{Auger}}
    = \bigl(C_n\,n + C_p\,p\bigr)\bigl(np - n_i^2\bigr),
    \label{eq:auger}
\end{equation}
where $C_n$ and $C_p$ are the electron and hole Auger capture coefficients.
The combined effective carrier lifetime accounting for both SRH and Auger
mechanisms is:
\begin{equation}
    \frac{1}{\tau_{n/p}}
    = \frac{1}{\tau_{\mathrm{Auger},n/p}}
    + \frac{1}{\tau_{\mathrm{SRH},n/p}},
    \label{eq:lifetime_combined}
\end{equation}
where $\tau_{\mathrm{Auger},n/p}$ is derived from \cref{eq:auger} evaluated
at the local carrier densities. \Cref{eq:lifetime_combined} defines the
effective recombination time used in the CHARGE simulation
of \cref{subsec:pd_charge}.

\emph{Radiative recombination} is modelled as:
\begin{equation}
    U_{\mathrm{rad}} = r_{\mathrm{opt}}\!\left(np - n_i^2\right),
    \label{eq:uradiative}
\end{equation}
where $r_{\mathrm{opt}}^{\mathrm{Ge}}\approx\SI{6.4e-14}{\centi\metre\cubed\per\second}$
at \SI{300}{\kelvin}, approximately four orders of magnitude smaller than for a
direct-gap III-V of comparable bandgap. In the fully depleted intrinsic region,
where $np\approx n_i^2$, the radiative rate vanishes identically and
\cref{eq:srh_dep} constitutes the sole generation mechanism. Simulated generation
and recombination rates confirm that SRH exceeds Auger by more than two orders of
magnitude and radiative recombination by more than five orders of magnitude across
$0\leq|V_{\mathrm{bias}}|\leq\SI{3}{\volt}$. This justifies the single-parameter
calibration strategy of \cref{subsec:pd_charge}, in which only $\tau_0$ is
adjusted to reproduce $I_d$.

Surface recombination at the Ge/SiO$_2$ and Ge/Si interfaces constitutes a
separate loss channel characterised by the surface recombination velocity $v_s$.
Its contribution to the dark current is proportional to the mesa perimeter $P_j$
rather than the junction area $A_j$, as formalised in
\cref{eq:jdark_surf,eq:idark_decomposed}. As device dimensions shrink to improve
bandwidth, the area-to-perimeter ratio $A_j/P_j$ decreases and surface
recombination becomes increasingly significant, setting a practical lower bound
on the achievable dark current density that surface passivation, for example
through Ge$_x$O$_y$ or Al$_2$O$_3$, must address.

\begin{figure}[H]
    \centering
    \missingfigure[figwidth=0.84\linewidth]{
        Three-panel concept figure. Panel (a): photon absorption in Ge
        creating an electron-hole pair. Panel (b): competing SRH, Auger, and
        radiative recombination pathways with their relative strengths in the
        depleted intrinsic region. Panel (c): qualitative reverse-bias trend
        showing SRH generation dominating the dark-current contribution in the
        operating range of the present device.}
    \caption{Carrier generation and recombination picture used throughout the
    electrical model. In the fully depleted intrinsic Ge layer, SRH generation
    remains the dominant dark-current mechanism, which justifies calibrating
    the device primarily through the lifetime parameter $\tau_0$.}
    \label{fig:generation_recombination}
\end{figure}

% ============================================================
\subsection[Design Synthesis: Linking Physics to Device Targets]%
           {Design Synthesis: Linking Physics to Device Targets}
\label{subsec:pd_synthesis}

The preceding sections have developed the physical foundations of photodetection
from first principles, and it is instructive at this point to consolidate the
key parameter interdependencies into a coherent engineering picture before
proceeding to the numerical device models of subsequent chapters.

\paragraph{The absorber-length/bandwidth Pareto front.}
For a waveguide-integrated detector, optical absorption and bandwidth are in
direct tension through the absorber length $L$. Responsivity benefits from
long $L$ via $\mathcal{A}(L) = 1 - e^{-\alpha L}$, while bandwidth degrades
primarily through the junction capacitance $C_j \propto A_j = W_j \cdot L$,
which enters the RC bandwidth \cref{eq:bw_rc}. Differentiating
$f_{3\mathrm{dB}}$ with respect to $L$ at fixed $W_i$ and $W_j$ yields the
responsivity-bandwidth product as the relevant composite figure of merit:
\begin{equation}
    \mathcal{R}\cdot f_{3\mathrm{dB}}
    \approx \frac{q\lambda}{hc}\!\left(1-e^{-\alpha L}\right)
    \cdot\left[\frac{1}{f_{\mathrm{tt}}^2}
    + \frac{(2\pi R_{\mathrm{tot}}\varepsilon_0\varepsilon_{\mathrm{Ge}}W_j L)^2}
           {W_i^2}\right]^{-1/2},
    \label{eq:R_BW_product}
\end{equation}
which is maximised by selecting $L$ such that the RC pole is pushed to roughly
the transit-time pole, i.e.\ $f_{\mathrm{RC}} \approx f_{\mathrm{tt}}$. For the
400GBASE-DR4 target requiring a per-lane bandwidth of \SI{53}{\giga\hertz} at
$\mathrm{PAM4}$ signalling, and a responsivity floor of
$\mathcal{R} \geq \SI{0.8}{\ampere\per\watt}$, \cref{eq:R_BW_product} constrains
the design space to approximately $L\in[\SI{10}{\micro\metre},\,\SI{20}{\micro\metre}]$
and $W_i\in[\SI{0.5}{\micro\metre},\,\SI{1.0}{\micro\metre}]$ for
$W_j = \SI{1}{\micro\metre}$ in the vertical PIN geometry.

\paragraph{Bias voltage and depletion completeness.}
The applied reverse bias $V_{\mathrm{bias}}$ must be large enough to fully
deplete the intrinsic layer and drive the carriers into velocity saturation.
The full-depletion voltage in the ideal PIN approximation is~\cite{SzeNg2006}:
\begin{equation}
    V_{\mathrm{FD}}
    = \frac{q\,N_{\mathrm{bg}}\,W_i^2}{2\varepsilon_0\varepsilon_{\mathrm{Ge}}} - V_{\mathrm{bi}},
    \label{eq:vfd}
\end{equation}
where $N_{\mathrm{bg}}$ is the unintentional background doping of the nominally
intrinsic layer, which is never zero in practice due to incorporation of boron or
phosphorus during epitaxial growth. For $N_{\mathrm{bg}} \approx
\SI{e16}{\per\centi\metre\cubed}$ and $W_i = \SI{1}{\micro\metre}$,
\cref{eq:vfd} gives $V_{\mathrm{FD}} \approx \SI{0.27}{\volt}$, well below
the standard operating voltage of \SI{1}{\volt}, confirming that full depletion
is achieved at low bias for realistic material quality. Above $V_{\mathrm{FD}}$,
the capacitance becomes approximately voltage-independent at the geometric value
of \cref{eq:cj_planar}, which simplifies RC bandwidth analysis. Operation below
$V_{\mathrm{FD}}$ should be avoided as it introduces a bias-dependent capacitance
and incomplete carrier collection, degrading both bandwidth and responsivity in a
non-linear, difficult-to-model manner.

\paragraph{Temperature dependence and thermal stability.}
The intrinsic carrier concentration of germanium follows
$n_i \propto T^{3/2}\exp(-E_g/2k_BT)$, where $E_g(T)$ itself decreases with
temperature according to the Varshni relation. Because $n_i^{\mathrm{Ge}}$ at
room temperature is already large, even modest self-heating in a high-power
optical link can substantially elevate the dark current through
\cref{eq:jdark_bulk}, degrading the SNR floor of \cref{eq:snr}. In a
\SI{400}{\giga\bit\per\second} DR4 transceiver module, thermal management
of the receiver die is therefore a system-level co-design requirement, not
merely a device-level footnote. The temperature coefficient of $I_d$ in
Ge-on-Si is typically measured to be in the range \SIrange{5}{8}{\percent\per\kelvin}
near room temperature, significantly steeper than silicon photodiodes, and this
sensitivity motivates both the careful thermal modelling performed alongside
the electrical characterisation of \cref{subsec:pd_charge} and the use of
bias-point tuning to compensate for operating-temperature drift in deployed
link hardware.
